\section{系统模型}

\subsection{参与方}

TrustAuditFlow 包含五类角色。

\textbf{数据拥有方。} 数据拥有方注册数据对象，定义审计策略，并生成初始承诺。
\textbf{链下存储节点。} 这些节点存储数据分片、校验分片或计算证据，并响应审计挑战。
\textbf{审计者。} 审计者抽样分片或证据对象，检查哈希或证明输出，并提交审计摘要。
\textbf{审计委员会。} 委员会从可信节点中选出，确认常规审计摘要、恢复状态和信任分更新。
\textbf{联盟链。} 联盟链存储对象状态、Merkle 根、哈希链锚点、恢复事件和节点信任状态。

\subsection{链上与链下边界}

系统遵循链下计算和链上锚定设计。链下节点执行分片、纠删码编码、本地存储、抽样校验、恢复和完整证据保留。链上记录紧凑证据：对象标识、编码参数、Merkle 根、批量审计摘要、恢复状态和信任分更新。链并不作为高频证明执行器，而是作为可信时间线、状态锚点和问责账本。

\subsection{对象与承诺}

设数据对象 $D$ 被编码为 $k$ 个数据分片和 $m$ 个校验分片。分片 $s_i$ 的哈希为 $h_i=H(s_i)$。系统基于所有分片哈希构建 Merkle 树，并得到根 $R_D$。链上对象状态存储为
\begin{equation}
  \langle id_D, k, m, R_D, policy \rangle,
\end{equation}
其中 $policy$ 包含抽样率、审计间隔、风险级别、委员会规模和恢复规则。

链下审计证据被组织为哈希链。对于审计事件 $e_t$，链状态为
\begin{equation}
  c_t = H(c_{t-1} \parallel e_t).
\end{equation}
系统周期性锚定 $c_t$ 或事件批次的 Merkle 根。只要锚定状态保持可用，后续对链下证据的删除或修改就能够被检测。

